\chapter{Reflection-positive rings}
\label{chap:rings}

To apply the spacetime-bootstrap criterion of Chapters~\ref{chap:criterion} and~\ref{chap:chain} to
a periodic spin chain, we must show that its partition functions satisfy the three
transfer-operator moment inequalities of Definition~\ref{def:moment-bounds}. For a nearest-neighbor
lattice Hamiltonian that is reflection positive across bonds~\cite{Osterwalder1973Axioms,
Osterwalder1975Axioms, Dyson1978Phase, Frohlich1978Phase, Biskup2009Reflection,
Jaffe2016Characterization}, spatial reflection plays the role of Euclidean time
reflection~\cite{Luscher1977Construction} and makes the column-to-column spatial transfer operator
self-adjoint~\cite{Suzuki1985Transfermatrix, TangTuWang}. In continuous imaginary time,
the spectral moment representation~\eqref{eq:moment-form} arises from a thermal matrix product
operator whose spatial transfer operator \(X_\beta\) is a compact self-adjoint operator on a
worldline path space, a construction that requires measure theory on Poisson configuration spaces
and the spectral theorem on infinite-dimensional Hilbert spaces. Instead, we discretize imaginary
time into finite brickwork Trotter slices~\cite{Suzuki1985Transfermatrix} (Definition~\ref{def:trotter})
so that each approximant is the trace of a power of a finite positive semidefinite spatial transfer
matrix (Theorem~\ref{thm:psd-transfer}); the resulting moment inequalities
(Lemma~\ref{lem:trotter-moment}) then pass to the continuous-time limit via the matrix
Lie--Trotter product formula (Theorem~\ref{thm:lie-trotter}) and
Lemma~\ref{lem:moment-bounds-limit} (Theorem~\ref{thm:rp-moment-bounds}).

\section{Ring Hamiltonians and bond reflection positivity}

Operators on an \(L\)-cell spin ring act on the tensor-product Hilbert space \((\C^d)^{\otimes L}\),
identified with complex wavefunctions on the product configuration space \(d^L\) for a finite local
basis \(d\) (Definition~\ref{def:ring-operators}). A reflection-positive ring
(Definition~\ref{def:rpring}) packages a nearest-neighbor bond interaction in the real factored
form~\eqref{eq:rp-bond}~\cite{Dyson1978Phase, Frohlich1978Phase} together with a real
symmetric on-site term \(o\) and an internal cell reflection \(F\), ensuring that the ring
Hamiltonian~\eqref{eq:rp-hamiltonian} is Hermitian (Lemma~\ref{lem:rpring-hermitian}) and that its
partition function \(\operatorname{Re}\Tr e^{-\beta H_L}\) (Definition~\ref{def:ring-spectrum})
equals the Boltzmann sum of Definition~\ref{def:partition-function} (Lemma~\ref{lem:ring-trace-exp}).

\begin{definition}[Local operators on a periodic chain]
  \label{def:ring-operators}
  \lean{CertifiedGaps.piTensor, CertifiedGaps.bondOp, CertifiedGaps.siteOp}
  \leanok
  For a finite local basis \(d\) and a periodic chain of \(L\) unit cells, operators act on complex
  functions on the product configuration space \(d^L\). The sitewise tensor product
  \(\bigotimes_i A_i\) has matrix entries \(\prod_i (A_i)_{\sigma_i \tau_i}\). A two-site operator
  \(B\) on cells \(j, k\) acts as \(B\) on \((\sigma_j, \sigma_k)\) and as the identity elsewhere;
  a single-site operator \(A\) at cell \(j\) is \(\bigotimes_i A_i\) with \(A_j = A\) and
  \(A_i = 1\) otherwise.
\end{definition}

\begin{definition}[Reflection-positive ring Hamiltonian]
  \label{def:rpring}
  \lean{CertifiedGaps.RPRing, CertifiedGaps.RPRing.bondMatrix, CertifiedGaps.RPRing.bond,
    CertifiedGaps.RPRing.site, CertifiedGaps.RPRing.hamiltonian}
  \leanok
  \uses{def:ring-operators}
  A reflection-positive ring with local basis \(d\) consists of real matrices
  \(C_\alpha \in \R^{d \times d}\) (\(\alpha \in \iota\)), each symmetric or antisymmetric, a real
  symmetric on-site operator \(o\), and an internal cell reflection \(F\) (a real symmetric
  involution satisfying \(F^2 = 1\) and \(F o F = o\), with \(F = 1\) by default). The
  nearest-neighbor bond interaction is
  \begin{equation}
    \label{eq:rp-bond}
    h = -\sum_\alpha C_\alpha \otimes F C_\alpha F,
  \end{equation}
  and the ring Hamiltonian on \(L\) cells is
  \begin{equation}
    \label{eq:rp-hamiltonian}
    H_L = \sum_{j=0}^{L-1} h_{j, j+1} + \sum_{j=0}^{L-1} o_j
  \end{equation}
  with \(j + 1\) taken cyclically modulo \(L\). A model with a multi-site unit cell is represented
  as a ring of cells in which \(d\) is the local cell basis and intra-cell interactions are
  absorbed into \(o\).
\end{definition}

\begin{remark}[Sufficient bond factorization versus the Jaffe--Janssens characterization]
  \label{rem:jaffe-janssens}
  Definition~\ref{def:rpring} formalizes the sufficient algebraic condition~\eqref{eq:rp-bond} for
  bond reflection positivity going back to Dyson, Lieb, and Simon~\cite{Dyson1978Phase} and
  Fr\"ohlich, Israel, Lieb, and Simon~\cite{Frohlich1978Phase}, extended by the internal cell
  involution \(F\) for multi-site unit cells. Every certified model in
  Chapters~\ref{chap:heisenberg}--\ref{chap:models} satisfies~\eqref{eq:rp-bond} directly via an
  explicit finite family of real symmetric or antisymmetric matrices \(C_\alpha\). The complete
  classification of reflection-positive spin and Majorana Hamiltonians by Jaffe and
  Janssens~\cite{Jaffe2016Characterization} establishes that, for any antiunitary reflection
  \(\Theta\), a bond interaction is reflection positive if and only if \(-h = \sum_\alpha C_\alpha
  \otimes \Theta(C_\alpha)\) after single-site shifts (equivalently, the bond realignment matrix is
  Hermitian positive semidefinite). The formalization uses
  only the sufficient decomposition~\eqref{eq:rp-bond} and does not require the converse
  classification direction of~\cite{Jaffe2016Characterization}.
\end{remark}

\begin{lemma}[Hermiticity of the ring Hamiltonian]
  \label{lem:rpring-hermitian}
  \lean{CertifiedGaps.RPRing.bondMatrix_isHermitian, CertifiedGaps.RPRing.oc_isHermitian,
    CertifiedGaps.RPRing.hamiltonian_isHermitian, CertifiedGaps.bondOp_isHermitian,
    CertifiedGaps.siteOp_isHermitian}
  \leanok
  \uses{def:rpring}
  The bond interaction \(h\) in~\eqref{eq:rp-bond}, the on-site operator \(o\), and the ring
  Hamiltonian \(H_L\) in~\eqref{eq:rp-hamiltonian} are Hermitian.
\end{lemma}
\begin{proof}
  \leanok
  Because the operators \(C_\alpha\) and the reflection \(F\) in Definition~\ref{def:rpring} are
  real, each summand of~\eqref{eq:rp-bond} satisfies
  \((C_\alpha \otimes F C_\alpha F)^\dagger = C_\alpha^{\mathsf T} \otimes F C_\alpha^{\mathsf T}
  F\),
  and the two minus signs from an antisymmetric factor \(C_\alpha\) cancel. Embedding Hermitian
  one- and two-site operators into the ring via Definition~\ref{def:ring-operators} preserves
  Hermiticity.
\end{proof}

\begin{definition}[Eigenvalue spectrum and partition function]
  \label{def:ring-spectrum}
  \lean{CertifiedGaps.RPRing.spectrum, CertifiedGaps.RPRing.partitionFunction}
  \leanok
  \uses{lem:rpring-hermitian}
  The spectrum \(E_L\) is the family of eigenvalues of \(H_L\) repeated according to multiplicity,
  and the partition function is \(Z_L(\beta) = \operatorname{Re}\Tr e^{-\beta H_L}\).
\end{definition}

\begin{lemma}[Spectral trace formula for the partition function]
  \label{lem:ring-trace-exp}
  \lean{CertifiedGaps.trace_exp_smul_eq_sum_of_isHermitian,
    CertifiedGaps.RPRing.partitionFunction_eq_sum}
  \leanok
  \uses{def:ring-spectrum, def:partition-function}
  For every Hermitian matrix \(A\) and real number \(c\), \(\Tr e^{cA} = \sum_i e^{c \lambda_i}\).
  In particular, \(\operatorname{Re}\Tr e^{-\beta H_L} = \sum_m e^{-\beta E_L(m)}\) agrees with the
  Boltzmann partition function of Definition~\ref{def:partition-function}.
\end{lemma}
\begin{proof}
  \leanok
  By the spectral theorem, \(A = U D U^\dagger\) for a unitary matrix \(U\) and a real diagonal
  eigenvalue matrix \(D\), so \(e^{cA} = U e^{cD} U^{-1}\) and the trace is invariant under
  conjugation.
\end{proof}

\section{Brickwork Trotter decomposition and the transfer matrix}

To exhibit the spectral moment representation~\eqref{eq:moment-form} at finite Trotter number \(M\),
we decompose each imaginary-time slice into a brickwork product~\eqref{eq:trotter-slice} of even
bonds, on-site terms, and odd bonds (Definition~\ref{def:trotter}). Expanding the linearized bond
factors (Lemma~\ref{lem:bond-expansion}) turns the \(M\)-slice trace on a ring of \(2N\) cells into
a sum, over imaginary-time bond configurations, of products of single-cell traces
\(\mathcal F_{pq}\) defined in~\eqref{eq:cell-transfer} (Definition~\ref{def:cell-transfer} and
Lemmas~\ref{lem:slice-sum}--\ref{lem:odd-word}); this sum contracts to
\(\Tr (\mathcal F^{\mathsf T} \mathcal F)^N\) in~\eqref{eq:transfer-factorization}
(Theorem~\ref{thm:transfer-factorization}). Reflection positivity guarantees that \(\mathcal F\) has
real entries (Lemma~\ref{lem:real-transfer}), so the two-cell spatial transfer matrix
\(T = \mathcal F^{\mathsf T} \mathcal F = \mathcal F^\dagger \mathcal F\) is positive semidefinite
(Theorem~\ref{thm:psd-transfer}).

\begin{definition}[Brickwork Trotter slice]
  \label{def:trotter}
  \lean{CertifiedGaps.RPRing.trotterSlice, CertifiedGaps.RPRing.trotterTrace,
    CertifiedGaps.evenBonds, CertifiedGaps.oddBonds}
  \leanok
  \uses{def:rpring}
  A first-order brickwork Trotter slice of step size \(\tau\) is the ordered product
  \begin{equation}
    \label{eq:trotter-slice}
    S_L(\tau) = \prod_{j \text{ even}} (1 - \tau h_{j,j+1}) \cdot \prod_i (1 - \tau o_i) \cdot
      \prod_{j \text{ odd}} (1 - \tau h_{j,j+1}),
  \end{equation}
  and the \(M\)-slice Trotter approximant of \(Z_L(\beta)\) is
  \(Z^{(M)}_L(\beta) = \operatorname{Re}\Tr S_L(\beta/M)^M\).
\end{definition}

\begin{lemma}[Rank-one decomposition of linearized bond factors]
  \label{lem:bond-expansion}
  \lean{CertifiedGaps.RPRing.one_sub_smul_bond, CertifiedGaps.RPRing.one_sub_smul_site}
  \leanok
  \uses{def:trotter}
  Let \(\tau \ge 0\), \(L \ge 2\), \(C_{\varnothing} = 1\), and \(C_\alpha^\tau = \sqrt\tau\,
  C_\alpha\). Then
  \(1 - \tau h_{j,j+1} = \sum_{\gamma \in \iota \sqcup \{\varnothing\}} (C^\tau_\gamma)_j \otimes
  (F C^\tau_\gamma F)_{j+1}\), and \(1 - \tau o_i\) is the single-site operator \(O = 1 - \tau o\)
  at cell \(i\).
\end{lemma}
\begin{proof}
  \leanok
  Substituting the bond decomposition~\eqref{eq:rp-bond} of \(-h\) from Definition~\ref{def:rpring}
  gives the sum over bond indices; the term \(\gamma = \varnothing\) is the identity operator
  because \(F^2 = 1\), and the remaining terms use \(\sqrt\tau \cdot \sqrt\tau = \tau\).
\end{proof}

\begin{lemma}[Trace of a matrix power as a sum over closed paths]
  \label{lem:closed-path}
  \lean{CertifiedGaps.trace_pow_eq_sum_closedPath, CertifiedGaps.piTensor_mul,
    CertifiedGaps.trace_piTensor, CertifiedGaps.list_prod_ofFn_sum}
  \leanok
  \uses{def:ring-operators}
  The trace of a matrix power expands as \(\Tr T^N = \sum_{q} \prod_{k} T_{q_k, q_{k+1}}\) over
  closed paths \(q = (q_0, \dots, q_{N-1}) \in D^N\) (with cyclic indices). Sitewise tensor
  products multiply factor by factor, \(\Tr \bigotimes_i A_i = \prod_i \Tr A_i\), and an ordered
  product of finite sums equals the sum of the ordered products.
\end{lemma}
\begin{proof}
  \leanok
  Induction on the exponent expands a matrix power as a sum over paths, and taking the trace
  restricts the sum to closed paths; the tensor-product and distributivity identities hold
  entrywise from Definition~\ref{def:ring-operators}.
\end{proof}

\begin{lemma}[Expansion of a Trotter slice into sitewise operators]
  \label{lem:slice-sum}
  \lean{CertifiedGaps.RPRing.trotterSlice_eq_sum, CertifiedGaps.RPRing.sliceSite_evenSite,
    CertifiedGaps.RPRing.sliceSite_oddSite}
  \leanok
  \uses{lem:bond-expansion, lem:closed-path}
  On a ring of \(2N\) cells with \(N \ge 1\), one Trotter slice~\eqref{eq:trotter-slice} expands as
  a sum over bond indices of sitewise tensor products. In each summand, even cell \(2k\) carries the
  local operator \(C^\tau_{b_k} O\, F C^\tau_{c_{k-1}} F\) and odd cell \(2k+1\) carries
  \(F C^\tau_{b_k} F\, O\, C^\tau_{c_k}\), where \(b_k\) indexes the even bond \((2k, 2k+1)\) and
  \(c_k\) indexes the odd bond \((2k+1, 2k+2)\).
\end{lemma}
\begin{proof}
  \leanok
  Expand each bond factor in~\eqref{eq:trotter-slice} by Lemma~\ref{lem:bond-expansion} and multiply
  the resulting sitewise tensor products by Lemma~\ref{lem:closed-path}, using the fact that each
  cell is incident to one even bond and one odd bond.
\end{proof}

\begin{definition}[Single-cell transfer matrix]
  \label{def:cell-transfer}
  \lean{CertifiedGaps.RPRing.timeWord, CertifiedGaps.RPRing.cellTransfer}
  \leanok
  \uses{lem:bond-expansion}
  For imaginary-time bond sequences \(p, q \in (\iota \sqcup \{\varnothing\})^M\) of length \(M\),
  the single-cell transfer matrix has entries
  \begin{equation}
    \label{eq:cell-transfer}
    \mathcal F_{pq} = \Tr \prod_{t=0}^{M-1} C^\tau_{p_t}\, O\, F C^\tau_{q_t} F,
  \end{equation}
  given by the trace of the time-ordered product of local operators on an even cell across the
  \(M\) Trotter slices.
\end{definition}

\begin{lemma}[Sublattice symmetry of the single-cell trace]
  \label{lem:odd-word}
  \lean{CertifiedGaps.RPRing.trace_oddWord}
  \leanok
  \uses{def:cell-transfer}
  On an odd cell, the trace of the time-ordered product satisfies
  \(\Tr \prod_{t=0}^{M-1} F C^\tau_{p_t} F\, O\, C^\tau_{q_t} = \mathcal F_{pq}\).
\end{lemma}
\begin{proof}
  \leanok
  Because \(F\) commutes with \(O\) and satisfies \(F^2 = 1\) (Definition~\ref{def:rpring}),
  conjugating each slice factor in~\eqref{eq:cell-transfer} on an even cell by the involution \(F\)
  gives the corresponding factor on an odd cell, and the trace is invariant under conjugation.
\end{proof}

\begin{theorem}[Spatial transfer-matrix factorization]
  \label{thm:transfer-factorization}
  \lean{CertifiedGaps.RPRing.trace_trotterSlice_pow_eq_sum,
    CertifiedGaps.RPRing.trace_trotterSlice_pow_eq_trace_transfer}
  \leanok
  \uses{lem:slice-sum, def:cell-transfer, lem:odd-word, lem:closed-path}
  For \(\tau \ge 0\), \(N \ge 1\), and every Trotter number \(M\),
  \begin{equation}
    \label{eq:transfer-factorization}
    \Tr S_{2N}(\tau)^M = \Tr \bigl(\mathcal F^{\mathsf T} \mathcal F\bigr)^N .
  \end{equation}
\end{theorem}
\begin{proof}
  \leanok
  Expanding the product of \(M\) slices~\eqref{eq:trotter-slice} by Lemmas~\ref{lem:closed-path}
  and~\ref{lem:slice-sum} expresses the ring trace as a sum over spacetime bond configurations of
  the product of single-cell traces. By Lemma~\ref{lem:odd-word}, both even and odd cells contribute
  entries of \(\mathcal F\) in~\eqref{eq:cell-transfer} indexed by the imaginary-time sequences on
  their two incident bonds. Summing over the even-bond sequences contracts each even cell with the
  adjacent odd cell into \(\mathcal F^{\mathsf T} \mathcal F\), and the remaining cyclic sum over
  the odd-bond sequences around the ring of \(N\) cell pairs equals
  \(\Tr (\mathcal F^{\mathsf T}\mathcal F)^N\) by Lemma~\ref{lem:closed-path}.
\end{proof}

\begin{lemma}[Reality of the single-cell transfer matrix]
  \label{lem:real-transfer}
  \lean{CertifiedGaps.RPRing.star_cellTransfer_apply,
    CertifiedGaps.RPRing.conjTranspose_cellTransfer}
  \leanok
  \uses{def:cell-transfer}
  The matrix \(\mathcal F\) in~\eqref{eq:cell-transfer} has real entries, so
  \(\mathcal F^\dagger = \mathcal F^{\mathsf T}\).
\end{lemma}
\begin{proof}
  \leanok
  Each entry \(\mathcal F_{pq}\) in~\eqref{eq:cell-transfer} is the trace of a product of real
  matrices, since Definition~\ref{def:rpring} expresses the negated bond interaction in
  \eqref{eq:rp-bond} as a sum of tensor products \(C_\alpha \otimes F C_\alpha F\) of real matrices
  and the on-site operator \(o\) is real.
\end{proof}

\begin{theorem}[Positive semidefiniteness of the two-cell transfer matrix]
  \label{thm:psd-transfer}
  \lean{CertifiedGaps.RPRing.exists_posSemidef_transfer}
  \leanok
  \uses{thm:transfer-factorization, lem:real-transfer}
  For every \(M\) and \(\tau \ge 0\) there is a positive semidefinite spatial transfer matrix \(T\),
  independent of the ring size, such that \(\Tr S_{2N}(\tau)^M = \Tr T^N\) for every \(N \ge 1\).
\end{theorem}
\begin{proof}
  \leanok
  Set \(T = \mathcal F^\dagger \mathcal F = \mathcal F^{\mathsf T} \mathcal F \succeq 0\) by
  Lemma~\ref{lem:real-transfer} and apply~\eqref{eq:transfer-factorization} from
  Theorem~\ref{thm:transfer-factorization}.
\end{proof}

\begin{lemma}[Spectral trace of a matrix power]
  \label{lem:trace-pow-eig}
  \lean{CertifiedGaps.trace_pow_eq_sum_eigenvalues}
  \leanok
  For every Hermitian matrix \(T\), \(\Tr T^N = \sum_i \lambda_i^N\).
\end{lemma}
\begin{proof}
  \leanok
  Diagonalizing the Hermitian matrix by the spectral theorem reduces the trace of its \(N\)-th
  power to the sum of the \(N\)-th powers of its eigenvalues.
\end{proof}

\begin{proposition}[Spectral moment representation from a positive semidefinite transfer matrix]
  \label{prop:moment-psd}
  \lean{CertifiedGaps.hasMomentForm_of_trace_pow}
  \leanok
  \uses{lem:trace-pow-eig, def:moment-form}
  If \(T \succeq 0\) and \(Z(2N) = \operatorname{Re}\Tr T^N\) for all \(N \ge 1\), then \(Z\)
  admits a spectral moment representation~\eqref{eq:moment-form} above length \(2\), with singular
  values \(\mu_i = \sqrt{\lambda_i(T)}\) (padded by zeros).
\end{proposition}
\begin{proof}
  \leanok
  By Lemma~\ref{lem:trace-pow-eig}, \(\operatorname{Re}\Tr T^N = \sum_i \lambda_i^N = \sum_i
  \mu_i^{2N}\)
  with \(\lambda_i \ge 0\), which has the form~\eqref{eq:moment-form} of
  Definition~\ref{def:moment-form}.
\end{proof}

\begin{lemma}[Spectral moment representation of Trotter approximants]
  \label{lem:trotter-moment}
  \lean{CertifiedGaps.RPRing.hasMomentForm_trotterTrace}
  \leanok
  \uses{thm:psd-transfer, prop:moment-psd}
  For every \(M\) and \(\beta > 0\), the Trotter partition function
  \(L \mapsto Z^{(M)}_L(\beta)\) of Definition~\ref{def:trotter} admits a spectral moment
  representation~\eqref{eq:moment-form} above length \(2\).
\end{lemma}
\begin{proof}
  \leanok
  Combine Theorem~\ref{thm:psd-transfer} at \(\tau = \beta/M\) with
  Proposition~\ref{prop:moment-psd}.
\end{proof}

\section{The Lie--Trotter product formula}

To pass from the \(M\)-slice Trotter approximants \(Z^{(M)}_L(\beta)\) (Definition~\ref{def:trotter})
to the exact partition function \(Z_L(\beta)\) (Definition~\ref{def:ring-spectrum}) at each fixed
ring length \(L\), we prove the matrix Lie--Trotter product formula (Theorem~\ref{thm:lie-trotter})
by telescoping powers in a normed algebra (Lemmas~\ref{lem:telescoping} and~\ref{lem:chernoff}).

\begin{lemma}[Telescoping bound for powers in a normed algebra]
  \label{lem:telescoping}
  \lean{CertifiedGaps.norm_pow_sub_pow_le}
  \leanok
  In a normed ring with \(\|1\| = 1\), if \(\|X\|, \|Y\| \le K\) and \(K \ge 1\), then
  \(\|X^k - Y^k\| \le k K^k \|X - Y\|\).
\end{lemma}
\begin{proof}
  \leanok
  Induct on \(k\) using the identity \(X^{k+1} - Y^{k+1} = X (X^k - Y^k) + (X - Y) Y^k\) and
  submultiplicativity of the norm.
\end{proof}

\begin{lemma}[First-order exponential product limit]
  \label{lem:chernoff}
  \lean{CertifiedGaps.tendsto_natCast_smul_sub_one,
  CertifiedGaps.tendsto_pow_of_tendsto_natCast_smul,
    CertifiedGaps.tendsto_pow_of_hasDerivAt}
  \leanok
  \uses{lem:telescoping}
  In a complete normed \(\C\)-algebra, if \(M (X_M - 1) \to A\) as \(M \to \infty\), then
  \(X_M^M \to e^A\). In particular, if \(g(0) = 1\) and \(g'(0) = A\), then
  \(g(1/M)^M \to e^A\).
\end{lemma}
\begin{proof}
  \leanok
  Compare \(X_M\) with \(Y_M = e^{A/M}\), whose \(M\)-th power is \(e^A\) and which also satisfies
  \(M(Y_M - 1) \to A\). With \(C = \|A\| + 1\), convergence implies that eventually
  \(\|X_M\|, \|Y_M\| \le 1 + C/M\), where \((1 + C/M)^M \le e^C\). Lemma~\ref{lem:telescoping} then
  bounds \(\|X_M^M - e^A\|\) by \(e^C \|M (X_M - Y_M)\|\), which tends to \(0\).
\end{proof}

\begin{theorem}[Matrix Lie--Trotter product formula]
  \label{thm:lie-trotter}
  \lean{CertifiedGaps.tendsto_list_prod_one_add_smul_pow,
    CertifiedGaps.hasDerivAt_list_prod_one_add_smul}
  \leanok
  \uses{lem:chernoff}
  For complex matrices \(h_1, \dots, h_k\) and \(c \in \R\),
  \begin{equation}
    \label{eq:lie-trotter}
    \Bigl(\prod_{i=1}^k \bigl(1 + \tfrac{c}{M} h_i\bigr)\Bigr)^M \longrightarrow
    \exp\Bigl(c \sum_{i=1}^k h_i\Bigr) \qquad (M \to \infty).
  \end{equation}
\end{theorem}
\begin{proof}
  \leanok
  The matrix polynomial \(\tau \mapsto \prod_i (1 + \tau h_i)\) has derivative \(\sum_i h_i\) at
  \(\tau = 0\); apply Lemma~\ref{lem:chernoff} in the matrix algebra equipped with the operator
  norm.
\end{proof}

\begin{lemma}[Convergence of Trotter partition functions]
  \label{lem:trotter-limit}
  \lean{CertifiedGaps.RPRing.tendsto_trotterTrace,
  CertifiedGaps.RPRing.sum_evenBonds_append_oddBonds}
  \leanok
  \uses{thm:lie-trotter, def:trotter, def:ring-spectrum}
  For every inverse temperature \(\beta\) and ring length \(L\),
  \(Z^{(M)}_L(\beta) \to Z_L(\beta)\) as \(M \to \infty\).
\end{lemma}
\begin{proof}
  \leanok
  The Trotter slice~\eqref{eq:trotter-slice} has the form \(\prod (1 + \tfrac{\beta}{M} h_i)\) over
  the list of negated even-bond, on-site, and odd-bond operators, whose sum is \(-H_L\) in
  \eqref{eq:rp-hamiltonian}. Apply~\eqref{eq:lie-trotter} from Theorem~\ref{thm:lie-trotter} and the
  continuity of the real trace \(A \mapsto \operatorname{Re}\Tr A\).
\end{proof}

\section{Moment bounds and the gap criterion for rings}

Combining the finite-\(M\) spectral moment representation (Lemma~\ref{lem:trotter-moment}) with the
Lie--Trotter limit (Lemma~\ref{lem:trotter-limit}) via Lemma~\ref{lem:moment-bounds-limit} shows
that every reflection-positive ring satisfies the transfer-operator moment inequalities of
Definition~\ref{def:moment-bounds} (Theorem~\ref{thm:rp-moment-bounds}). An admissible rational
doubling certificate (Definition~\ref{def:chain-certificate}) and two finite-size seed bounds then
immediately imply a uniform spectral gap (Theorem~\ref{thm:rp-gap}).

\begin{theorem}[Transfer-operator moment inequalities for reflection-positive rings]
  \label{thm:rp-moment-bounds}
  \lean{CertifiedGaps.RPRing.spatialMomentBounds, CertifiedGaps.RPRing.hasSpatialMomentBounds}
  \leanok
  \uses{lem:trotter-moment, lem:trotter-limit, lem:moment-bounds-limit, lem:ring-trace-exp}
  For every reflection-positive ring (Definition~\ref{def:rpring}) and every \(\beta > 0\), the
  partition function \(L \mapsto Z_L(\beta)\) satisfies the transfer-operator moment inequalities
  (1)--(3) of Definition~\ref{def:moment-bounds} above length \(2\). Equivalently, the eigenvalue
  families \(E_L\) satisfy them.
\end{theorem}
\begin{proof}
  \leanok
  The Trotter approximants admit a spectral moment representation~\eqref{eq:moment-form} by
  Lemma~\ref{lem:trotter-moment} and converge at each fixed length by Lemma~\ref{lem:trotter-limit},
  so Lemma~\ref{lem:moment-bounds-limit} gives the moment inequalities of
  Definition~\ref{def:moment-bounds} for the matrix traces. Rewriting the traces as eigenvalue sums
  via Lemma~\ref{lem:ring-trace-exp} yields the spectral statement.
\end{proof}

\begin{remark}
  Passing to the continuous-time limit in Theorem~\ref{thm:rp-moment-bounds} does not require a
  Trotter error bound that is uniform in \(L\), because the limit \(M \to \infty\) is taken at each
  fixed ring length inside the three moment inequalities of Definition~\ref{def:moment-bounds}
  before the scale-doubling recurrence is iterated. The transfer-operator moment inequalities hold
  unconditionally for every reflection-positive ring, including gapless and symmetry-broken chains;
  whether the model is gapped is decided entirely by the two finite-size seed inequalities in
  Theorem~\ref{thm:rp-gap}.
\end{remark}

\begin{definition}[Spatial and thermal purities of a ring]
  \label{def:rp-purities}
  \lean{CertifiedGaps.RPRing.spatialPurity, CertifiedGaps.RPRing.thermalPurity}
  \leanok
  \uses{def:ring-spectrum}
  The ring purities \(\Ps(n, \beta)\) and \(\Pt(L, \beta)\) are defined as in
  Definition~\ref{def:purities}, computed from \(Z_L(\beta) = \operatorname{Re}\Tr e^{-\beta H_L}\).
\end{definition}

\begin{theorem}[Spacetime-bootstrap criterion for a reflection-positive ring]
  \label{thm:rp-gap}
  \lean{CertifiedGaps.RPRing.gapGT_of_chainCertificate}
  \leanok
  \uses{thm:rp-moment-bounds, thm:gap-of-chain-certificate, def:rp-purities, lem:ring-trace-exp}
  Let \(n_0 \ge 2\) be even, \(\beta_0 > 0\), and \(c\) an admissible rational doubling certificate
  (Definition~\ref{def:chain-certificate}) with initial defects \((s_0, t_0)\) and decay base \(r\).
  If
  \begin{equation}
    \label{eq:rp-seed-hyp}
    \Ps(n_0, \beta_0) \ge 1 - s_0, \qquad \Pt(2n_0, \beta_0) \ge 1 - t_0,
  \end{equation}
  then every even ring of \(L \ge n_0\) cells has a unique ground state and a spectral gap greater
  than \(\log r / \beta_0\).
\end{theorem}
\begin{proof}
  \leanok
  Apply Theorem~\ref{thm:gap-of-chain-certificate} with the moment inequalities of
  Theorem~\ref{thm:rp-moment-bounds}, rewriting the purities~\eqref{eq:rp-seed-hyp} by
  Lemma~\ref{lem:ring-trace-exp}.
\end{proof}

\section{Sitewise unitary transformations and sublattice rotations}

Antiferromagnetic spin chains are often reflection positive only after a local basis change, such
as applying the \(\pi\)-rotation \(e^{i\pi S^y}\) on one sublattice of a bipartite
lattice~\cite{Dyson1978Phase, Biskup2009Reflection}. Conjugating the ring Hamiltonian
\eqref{eq:rp-hamiltonian} by a sitewise tensor product of local unitaries
(Lemma~\ref{lem:conj-ops} and Theorem~\ref{thm:conj-hamiltonian}) preserves the full eigenvalue
spectrum (Theorem~\ref{thm:conj-spectrum}), so a spectral gap established in the reflection-positive
frame holds identically in the physical laboratory frame.

\begin{lemma}[Action of sitewise unitaries on local operators]
  \label{lem:conj-ops}
  \lean{CertifiedGaps.piTensor_conj, CertifiedGaps.conj_bondOp_kronecker,
    CertifiedGaps.conj_siteOp, CertifiedGaps.piTensor_conjTranspose_mul_self}
  \leanok
  \uses{def:ring-operators}
  Let \(U = \bigotimes_i a_i\) with \(a_i a_i^\dagger = 1\). Then \(U\) is unitary,
  \(U (A \otimes B)_{jk} U^\dagger = (a_j A a_j^\dagger \otimes a_k B a_k^\dagger)_{jk}\) for
  \(j \ne k\), and \(U A_j U^\dagger = (a_j A a_j^\dagger)_j\).
\end{lemma}
\begin{proof}
  \leanok
  Sitewise tensor products in Definition~\ref{def:ring-operators} multiply factor by factor, and
  \(a_i a_i^\dagger = 1\) cancels on every cell where the operator acts as the identity.
\end{proof}

\begin{theorem}[Sitewise unitary conjugation of the ring Hamiltonian]
  \label{thm:conj-hamiltonian}
  \lean{CertifiedGaps.RPRing.conj_hamiltonian}
  \leanok
  \uses{lem:conj-ops, def:rpring}
  For \(L \ne 1\),
  \begin{equation}
    \label{eq:conj-hamiltonian}
    U H_L U^\dagger = -\sum_j \sum_\alpha (a_j C_\alpha a_j^\dagger)_j \otimes
    (a_{j+1} F C_\alpha F a_{j+1}^\dagger)_{j+1} + \sum_j (a_j o a_j^\dagger)_j.
  \end{equation}
\end{theorem}
\begin{proof}
  \leanok
  Conjugating each bond term~\eqref{eq:rp-bond} and on-site term in~\eqref{eq:rp-hamiltonian} by
  Lemma~\ref{lem:conj-ops} gives~\eqref{eq:conj-hamiltonian}, since adjacent cells satisfy
  \(j \ne j+1\) whenever \(L \ge 2\).
\end{proof}

\begin{theorem}[Unitary invariance of the eigenvalue spectrum]
  \label{thm:conj-spectrum}
  \lean{CertifiedGaps.eigenvalues_eq_of_conj, CertifiedGaps.RPRing.eigenvalues_eq_spectrum}
  \leanok
  \uses{thm:conj-hamiltonian, def:ring-spectrum}
  Unitary conjugation preserves eigenvalues with multiplicity. Hence any Hermitian operator \(H\)
  satisfying \(H = U H_L U^\dagger\) shares the eigenvalue spectrum \(E_L\) of the ring
  (Definition~\ref{def:ring-spectrum}).
\end{theorem}
\begin{proof}
  \leanok
  Because \(\det(t - U A U^\dagger) = \det(t - A U^\dagger U) = \det(t - A)\), the characteristic
  polynomial is invariant under unitary conjugation, and the eigenvalue family of a Hermitian
  matrix is determined by its characteristic polynomial.
\end{proof}

